\section{Discussion and Limitations}
\label{sec:discussion}

\subsection{Security Envelope and Trade-offs}

The fallback trades cross-chain withdrawals for in-chain liveness. We
discuss the trade along three axes: liveness inside the chain, safety of
cross-chain assets, and the cost of griefing.

\textbf{Liveness inside the chain.} In \textsc{Sovereign}, block production,
in-chain transactions, application logic, and the forced-transaction queue
continue. Cross-chain withdrawals stop. A hard-gated chain under the same
degradation commits only a fraction of its blocks and does not resume after
the layer returns (Section~\ref{subsec:rq2}).

\textbf{Safety of cross-chain assets.} The prototype has no bridge and no
real assets: the application mounts neither \texttt{x/auth} nor
\texttt{x/bank}, transactions are unsigned envelopes, and a withdrawal is
recognized by the substring \texttt{TX\_WITHDRAWAL} in the transaction bytes.
The circuit breaker is therefore evaluated as a policy: it shows that the
decision to admit or reject withdrawals is agreed by all validators and
applied identically, not that any economic loss is prevented.

\textbf{Cost of griefing.} An adversary that degrades the peripheral view of
\emph{all} honest validators -- for example by delaying their Bitcoin
connectivity -- can drive the chain into \textsc{Sovereign} and freeze
withdrawals for as long as the degradation lasts (P4,
Section~\ref{subsec:rq1}). This is a new attack surface relative to a chain
that ignores its periphery, but not relative to a hard-gated chain: the same
attacker halts the hard-gated chain entirely, whereas Engram freezes only
withdrawals. Lemma~\ref{lem:honest-range} prevents Byzantine validators from causing the freeze
alone: honest validators whose peripheral view is degraded must supply the
rest of the majority. Together with the Byzantine validators they must hold
more than half of the certificate weight, and the proposer's choice of quorum
subset sets that weight. At $n{=}7$ with two colluders, an honest proposer that
includes all seven observations requires two degraded honest validators, and a
Byzantine proposer that selects a minimal quorum of five requires one.
In a design without a quorum-backed health input, one node's view
suffices. The freeze is the intended response to an
unverifiable periphery, not a safety violation. Its economic cost to users and
the attacker's cost of sustaining it are not quantified.

\textbf{Repeated triggering of \textsc{Suspicious}.} Backoff on
\textsc{Recovering} already raises the price of forcing repeated failed
recoveries. Without the adaptive threshold of Section~\ref{subsec:rq3}, the
\textsc{Anchored}$\leftrightarrow$\textsc{Suspicious} boundary has no such
protection beyond the bounded forced-transaction queue; the adaptive
threshold is meant to close this asymmetry. In the simulation of
Section~\ref{subsec:rq3} it absorbs a naive repeated adversary (1 demotion in
1200 blocks against 300 for a static threshold of 2) but gains little over a
static threshold of 4, and the live sweep suppressed flapping under jitter no
better than a static threshold of 2; the live repeated-trigger runs are
\pending{report the live repeated-trigger runs of the adaptive threshold.}

\textbf{Hysteresis: noise against concealment.} A higher down threshold $k$
tolerates longer bursts of noise, and it also tolerates a longer run of reports
that a Byzantine proposer interleaves with honest ones. Within the assumption
the choice of $k$ does not matter: for each of the 21 placements of two
Byzantine proposers in the seven-slot rotation, a real DA outage was detected
after exactly $k$ blocks ($k\in\{2,3,4\}$, simulation). With three Byzantine
validators, outside the assumption, a Byzantine proposer can keep a certificate
that reports a healthy DA layer, and an honest proposer then reports the outage
in the blocks it proposes; the outage is detected only if enough honest
proposers follow one another to reach $k$. In the simulation it was detected for all 35
placements at $k{=}2$, missed for 14 at $k{=}3$ and for 28 at $k{=}4$, and with
four Byzantine validators it was missed for every placement. The live check at
$k{=}2$ agrees: with three Byzantine validators the outage was detected in 5 of
5 runs. The adaptive threshold exposes the same
weakness only while its flap level is high (0 of 35 placements hidden at level 0,
28 at level 1, all 35 at level 2), whereas a static threshold of 4 or more is
exposed permanently. A rule that demotes when at least $m$ of the last $w$ blocks are bad,
instead of after a consecutive run, is not cut by interleaved reports. In a
reduced model of the \textsc{Anchored}$\leftrightarrow$\textsc{Suspicious} edge,
validated against the keeper-based simulation above (static thresholds and the
adaptive $2\to4$ reproduce its flaps, detection delays and hidden placements),
4 of 8 hid the outage in 0 of 35 placements, against 28 for a static threshold
of 4, at 0.8 flaps against 0.2 and the same 4-block detection. It does not
remove the cost of being fast: its detection is as slow as the static
threshold, and an adversary that sends bursts of $m$ blocks still obtains 200
demotions in 1200 blocks, where the adaptive threshold allows 1. Combining the
two (an adaptive window) gave 4.6 flaps and 172 demotions, so no single rule
won on all four criteria. We did not implement the window rule. The live check at $k{=}4$ (three Byzantine validators with biased certificates, all seven bridges stopped for 120\,s, a fresh chain per run)
agrees with the simulation for adjacent placements and not for interleaved ones. With three consecutive Byzantine slots
the outage was detected in 3 of 3 runs, as predicted. With slots 0, 2 and 4 of the rotation, which the simulation predicts
would hide it, the chain also left \textsc{Anchored} and reached \textsc{Sovereign} in 3 of 3 runs, with 6 to 9 blocks outside the honest range
per run (6 to 12 over all six). We have not isolated why the prediction failed for the interleaved placement, so the concealment at $k{=}4$ rests on the simulation alone.

\textbf{Proposer bias within the honest range.} When honest observations
straddle a threshold, a proposer can choose which honest value the
certificate selects (Section~\ref{subsec:health-certificate}). This cannot
produce a transition that no honest view supports, but a sequence of
Byzantine proposers can bias the timing of transitions. A persistent split
straddling an FSM threshold (e.g.\ three versus four validators on each side)
has not been tested live; the RQ4 scan shows no extra rounds in the unfaulted case, and the
split case itself remains untested.

\textbf{Restart after a long chain.} The publisher's contiguous height is
kept in memory. A validator that restarts therefore republishes the chain from
height 0 and reports DA as failed until it has caught up. In one run, four
validators restarted on a chain of about 7{,}400 heights; the certified DA
health stayed failed and the chain remained in \textsc{Sovereign} for more than
half an hour. Persisting the contiguous height would remove the dependence on
chain length; the prototype does not.

\textbf{History divergence at re-anchoring.} During \textsc{Sovereign}, the
ordering of blocks rests on the validator quorum alone. A coalition
exceeding $f$ could produce a \textsc{Sovereign} history that differs from
one it would have produced under Bitcoin timestamping; the recovery proof
certifies the history's internal legality and quorum signatures, not its
uniqueness beyond the $f$ bound. The withdrawal lock limits what such a
history can extract.

\textbf{Data availability of the \textsc{Sovereign} interval.} During \textsc{Sovereign}, the interval's block data is held by the validators' own block stores, so its
recoverability depends on those validators until it is published. The
backfill gate of Section~\ref{subsec:reanchoring} makes publication a
precondition for returning to \textsc{Anchored}.

\textbf{The unproven tail: liveness against trust.} A recovery proof is built
while the chain keeps producing blocks, so by the time it is submitted the
tip stands beyond the last header it covers. The gate accepts the proof only
if that unproven tail is at most a bound, \texttt{MaxUnprovenTailBlocks}, so
the bound must be at least the proving latency times the block rate. The bound
sets a trade-off. A small bound keeps almost all of the \textsc{Sovereign}
interval under proof when withdrawals reopen, but a proof that takes longer
than the bound allows to build never satisfies the gate, and the chain stays in
\textsc{Recovering}. A larger bound lets slower proofs through, at the price
that the last blocks before the chain returns to \textsc{Anchored} are
certified by the validator quorum alone; a coalition above $f$ could rewrite
that stretch without the proof detecting it, which is the same residual
exposure as in the history-divergence paragraph above, now bounded by the
tail. At about 0.8 blocks per second, the default bound of 128 blocks is about
2.7 minutes and a bound of 512 about 10.7 minutes.

\textbf{Measured proving latency.} In the first live deployment with chunk
aggregation, the prover built a 512-header aggregate from two chunks in about
4.3 minutes and a 514-header aggregate from three in about 6.6 minutes; the
tip then stood 211 and 321 blocks beyond the last covered header, close to the
product of latency and block rate. Both proofs were accepted by
\texttt{SubmitRecoveryProof}, but with the bound at 128 neither made the
recovery flag hold, and the chain stayed in \textsc{Recovering}. A single
256-header proof takes about two minutes, about 110 blocks, and fits, but
with little margin: at a block rate about 17\% higher, or a slower prover, it
would not, so the recovery times of the proof-only runs hold for this
testbed's block rate and proving speed. With the bound raised to 512 blocks and
everything else unchanged, the chain left \textsc{Recovering} through an
aggregate proof in both of two deployments. In the first the prover built a
514-header proof from three chunks in 6.5 minutes, and the tip stood 309
blocks (about 6.4 minutes) beyond the last covered header when it was
submitted; in the second a 512-header proof from two chunks took 4.6 minutes,
with the tip 224 blocks (about 4.7 minutes) beyond. All seven validators were
then in \textsc{Anchored}, 1{,}023 and 888 seconds after the chain started.
Proof-only recovery was on, so only the proof could have restored it. The
interval proved in these runs is the one since genesis, not one following an
outage injected into a healthy chain, and the runs with the bound at 128 are a
single deployment.

Our prover builds the chunk proofs one after another, so latency grows with
the number of chunks. The chunk proofs are independent given the chunk
boundary roots, so they could be built on separate machines and joined by the
aggregator circuit. That would make the latency roughly that of one chunk plus
one aggregation, about three minutes on our testbed by our estimate from the
single-chunk and aggregation timings, which still exceeds the default bound of
128 blocks. Going below it would also need faster proving of each chunk, for
example on dedicated multi-core machines or with hardware acceleration such as
GPUs. The aggregation step still needs all chunk proofs on one machine. We did
not measure distributed proving or hardware acceleration. The validators
verify the submitted proof during block execution, so the verification cost
does not change with how the proof was built.

\subsection{Interpretation of Results}

The formal and empirical results cover complementary ground. The TLA+ results
establish, per layer and within bounded configurations, that the FSM's safety
invariants hold and that removing any of four defensive mechanisms reproduces
the violation it targets; Lemma~\ref{lem:honest-range} bounds the certified
health input by honest observations. The live evaluation shows that the
composed implementation behaves as specified under real peripheral failures.
Because the cross-layer refinement is open, the whole-system claim rests on
the empirical side. Single-trial live results (composite fault, seven-scenario
sweep, P4, reorg depth) show that correct behavior occurs without
characterizing its distribution; the repeated experiments (hysteresis sweeps,
the five-repeat overhead benchmark, the 2{,}000-trial detector study) do
characterize distributions.

\textbf{Unexpected results.} Three results need a causal account.
\begin{itemize}
  \item \emph{Block rate.} Against a single local vanilla process, Engram's
  block interval is longer, which at first suggested a cost of the extended
  object. Seven vanilla validators on the same hosts give the same interval,
  with vote extensions on or off, so the gap is the cost of the
  seven-validator network and Engram's own addition is about \resOverheadPct{}
  (approximate 95\% interval \resOverheadCI{}; Section~\ref{subsec:rq4}).
  \item \emph{Hysteresis threshold.} A static threshold of 2 flapped under
  wide-area jitter (9.0 flaps per validator), and thresholds of 4, 6 and 8 did
  not (Section~\ref{subsec:rq3}; three runs for 2 and 6, one for 4 and 8). We
  did not isolate which sensor crosses its threshold; a run with the jitter
  profile removed, and per-predicate logging, would show it.
  \item \emph{Reorg depth.} Reorgs of depth 3--5 orphaned confirmed blocks
  without a reaction, while depth 15 escalated through the
  suspicious-duration bound (Appendix~\ref{app:reorg}). Whether the
  anchor-reverification flag was ever set is not established, because the
  per-validator sensor value was not captured; logging it and sweeping depths
  6--14 would locate the boundary.
\end{itemize}

\subsection{Limitations}
\label{subsec:limitations}

\begin{itemize}
  \item \emph{Testbed.} All results come from one host with accelerated
  Bitcoin and single-bridge DA trust; the backfill gate rests on a
  quorum-reported contiguous height, not a proof of retrievability
  (Section~\ref{subsec:sim-to-real}).
  \item \emph{Proof gate.} The recovery-validity flag holds if a verified
  proof was submitted within the permitted unproven tail of the tracked tip,
  or if a heuristic holds: the last submitted Bitcoin checkpoint has reached
  the anchored depth. A protocol parameter disables the heuristic, so that
  only a verified proof satisfies the gate. The RQ2 trials and the five repeated composite runs used proof-only
  recovery with an unproven-tail bound of 128 blocks; the first composite run,
  the seven-scenario run and the data-availability runs of
  Section~\ref{subsec:rq3} were made with the heuristic enabled, so their
  time in \textsc{Recovering} does not isolate the cost of the recovery proof.
  Recovery through a proof also depends on the unproven-tail bound and on the
  testbed's block rate and proving speed, as discussed under the unproven tail
  above.
  \item \emph{Byzantine scale.} The single-actor attacks involve one
  Byzantine actor at a time; the $f{=}2$ collusion suite
  (Section~\ref{subsec:rq1}) addresses colluding validators.
  \item \emph{Fixed validator set.} The prototype assumes one validator set
  throughout a degraded interval (Section~\ref{subsec:threat-model}).
  \item \emph{Formal scope.} The cross-layer refinement is open, and the
  health certificate is modelled with Boolean observations in the server
  layer; numeric aggregation is covered by Lemma~\ref{lem:honest-range} only
  (Section~\ref{sec:formal}).
  \item \emph{Load range.} Load was tested up to 100 tx/s; behavior at
  higher rates is not characterized.
\end{itemize}

\subsection{Threats to Validity}

\textbf{Internal validity.} Host sharing can confound cost measurements.
Repeated measurements with confidence intervals, and a vanilla comparison on
the same host under the same load, mitigate this; the host still affects
absolute values.

\textbf{External validity.} Results at seven validators on one host need not
extrapolate to larger or geographically distributed validator sets; they are
evidence of correct behavior and of the source of the overhead, not production
performance figures.

\textbf{Construct validity.} AppHash agreement and height progression are the
conventional operational proxies for safety and liveness in CometBFT-family
systems. AppHash agreement alone cannot detect an illegal transition agreed on
by all validators; the semantic replay check of Section~\ref{subsec:rq1}
addresses this gap.

\subsection{Future Work}

Three directions remain beyond the items marked in this paper. First,
closing the cross-layer refinement obligation, and extending the
certificate model from Boolean to numeric observations, would extend the
formal results to the composed system; deductive methods~\cite{padon2016ivy,lotan2025implicit}
are a natural complement where bounded model checking does not scale.
Second, replacing the fixed-arity aggregator with incremental recursion
would make the recovery artifact constant-size for intervals of any length,
and proving the chunks in parallel or on dedicated hardware would shorten the
recovery whose latency the unproven-tail discussion measures. Third,
replacing single-bridge DA trust with sampling across independent peers
would let the DA sensor detect partial withholding rather than only outages.
A further step is the external verifier named in Section~\ref{sec:protocol}: a
light client or bridge on another chain that accepts a re-anchoring checkpoint
only with a valid proof. We have not built one or measured its verification
cost.
